\documentclass[aps,prl,twocolumn,floatfix,superscriptaddress,nofootinbib]{revtex4-2}

\usepackage{amsmath,amssymb,bm}
\usepackage{graphicx}
\usepackage{xcolor}
\usepackage{hyperref}

\hypersetup{
  colorlinks=true,
  linkcolor=blue,
  citecolor=blue,
  urlcolor=blue
}

\newcommand{\Tr}{\operatorname{Tr}}
\newcommand{\avg}[1]{\overline{#1}}
\newcommand{\Gtraj}{G_{\xi}}
\newcommand{\Gbar}{\overline{G}}
\newcommand{\Gps}{G_{\rm ps}}
\newcommand{\calL}{\mathcal{L}}
\newcommand{\calC}{\mathcal{C}}
\newcommand{\calD}{\mathcal{D}}
\newcommand{\calE}{\mathcal{E}}
\newcommand{\calN}{\mathcal{N}}
\newcommand{\dd}{\mathrm{d}}
\newcommand{\nsh}{n_{\rm shell}}
\newcommand{\atriv}{\alpha_{\rm triv}}
\newcommand{\atopo}{\alpha_{\rm top}}
\newcommand{\Dlam}{\Delta_{\lambda}}
\newcommand{\Dp}{\Delta_{\rm p}}

% ---- figure placeholder helper ----
\newcommand{\figplaceholder}[2]{%
  \fbox{\begin{minipage}[c][#1][c]{0.96\linewidth}
    \centering\small\textsf{#2}
  \end{minipage}}}

\begin{document}

\title{Steering fermions to chiral criticality}

\author{Asadullah Bhuiyan}
\affiliation{Department of Physics, Cornell University, Ithaca, New York 14853, USA}
\author{Haining Pan}
\affiliation{Department of Physics and Astronomy, Center for Materials Theory, Rutgers University, Piscataway, New Jersey 08854, USA}
\author{Chao-Ming Jian}
\affiliation{Department of Physics, Cornell University, Ithaca, New York 14853, USA}

\date{\today}

\begin{abstract}
We study the boundary of a topological monitored dynamics using a postselection-free adaptive Gaussian circuit in two dimensions.  A spatially patterned schedule of finite-support measurements and local feedback stabilizes a Chern region with a programmed domain wall.  The resulting pure, charge-fluctuating trajectory ensemble has log-chord entanglement consistent with a $c=1$ interface mode and controlled-window correlations consistent with $\Delta=2$.  Oppositely signed modular-charge drift on the two walls provides evidence for chirality, while a wall-localized near-marginal tangent sector connects the critical states to the dynamics that prepares them.  The trajectory-averaged state is instead mixed, showing that the universal boundary physics belongs to the monitored ensemble rather than to its average.
\end{abstract}

\maketitle

\emph{Introduction.---}One of the defining features of a topological phase is the bulk--boundary correspondence: an interface between bulks with distinct topological invariants supports modes whose existence does not follow from any local order parameter~\cite{CallanHarvey1985,Hatsugai1993,HasanKaneRMP,QiZhangRMP,ChiuTeoSchnyderRyu2016}.  For a Chern insulator, this gives a chiral fermion confined to the interface.  The corresponding question is less settled when the bulk is not governed by a Hamiltonian or a unitary evolution, but by a random sequence of measurements and feedback.  The problem is especially sharp at an internal domain wall, where the dynamics itself changes topology in space.  In that setting, one must characterize both the topology of the dynamics and the quantum states that its boundary prepares.

Adaptive circuits, built from measurements followed by outcome-dependent operations, provide a natural setting for this question~\cite{Skinner2019,LiChenFisher2018,Fisher2023,GullansHuse2020,Roy2020,Puente2024,Smith2024}.  Following Ref.~\cite{BhuiyanPanJian2026}, we call such a phase a \emph{topological monitored quantum dynamics} when a bulk invariant of the complete physical instrument ensemble remains well defined under local, symmetry-preserving spacetime disorder and is inherited by its steady-state trajectories.  The complete ensemble is essential.  A Floquet phase is defined by a repeatable unitary and its quasienergy spectrum~\cite{Kitagawa2010,Lindner2011,Rudner2013,Harper2020}, whereas a monitored circuit realizes a stochastic sequence of generally noninvertible operations.  A non-Hermitian description selects a no-click or otherwise postselected branch and omits the complementary outcomes and their Born probabilities~\cite{KellsMeidanRomito2023}.  The unconditional channel instead discards the record.  It may admit a Lindblad description in an appropriate weak-update limit, but a discrete averaged channel is not in general generated by a time-independent Lindbladian.  Moreover, under the assumptions of strict finite range, purity, and finite-rate relaxation, a local Lindbladian cannot prepare a unique pure Chern state in two dimensions~\cite{Diehl2011,Bardyn2013,BudichZollerDiehl2015,Goldstein2019}.

Earlier work established a symmetry-based classification of monitored free-fermion dynamics and related its topology to localization and state preparation~\cite{BhuiyanPanJian2026,JianBauerKeselmanLudwig2022,Ludwig2013,XiaoKawabata2024,KellsMeidanRomito2023,NahumSkinner2020}.  Xiao and Kawabata showed that an open boundary of a topological monitored dynamics produces gapless modes in the Lyapunov spectrum of the stochastic transfer matrix~\cite{XiaoKawabata2024}.  This identifies a boundary of the dynamical ensemble, but does not by itself specify the observables or universality of a physically prepared state.  Oshima \emph{et al.} characterized topological Lyapunov sectors at measurement-induced transitions in one-dimensional systems~\cite{Oshima2025}, while Pan, Shapourian, and Jian obtained Majorana zero modes at domain walls of genuinely monitored chains without postselection~\cite{PanShapourianJian2025}.  In two spatial dimensions, Ref.~\cite{BhuiyanPanJian2026} constructed a class-A adaptive protocol and observed anomalously slow purification near its topological domain walls.  Whether the surviving wall sector is critical and chiral, and how it appears in prepared-state observables, remained open.

Postselection changes this state-preparation question.  By conditioning on an appropriately chosen sequence of successful measurement outcomes, one can formally reproduce imaginary-time filtering $e^{-\beta H}$ toward a prescribed ground state, generally with a success probability that becomes exponentially small~\cite{LiZou2023,GarrattAltman2024}.  The protocol studied here does not condition on a rare record: every measurement outcome is sampled according to the Born rule and retained.

The circuit implements its target through localized, nonorthogonal overcomplete-Wannier modes constructed from the bands of a Chern insulator~\cite{BhuiyanPanJian2026}.  It measures the occupation of each mode and assigns the upper-band modes to be empty and the lower-band modes to be filled.  When a Born-sampled outcome disagrees with its target, a record-conditioned fermionic SWAP with an ancillary layer removes or injects a particle.  The enlarged system evolves unitarily during this exchange, while the induced operation on the physical layer is the required gain or loss.  Truncating the measured modes gives a strictly finite-support protocol.  A domain wall can then be programmed either by using modes associated with distinct bulk topologies on complementary regions or by terminating the support of a uniformly topological controller at a chosen boundary.  Fermion counting is therefore both the microscopic implementation and the observable record of the controller, rather than a separate counting field imposed on an otherwise fixed dynamics.

The object prepared at the wall is not a single target wavefunction.  It is an ensemble of pure Gaussian trajectories whose microscopic occupations and total charge fluctuate from record to record, while the bulk topology and long-distance interface physics are shared.  Averaging over those records produces a different, mixed state.  This distinction matters because entanglement and squared correlations are nonlinear in the trajectory covariance: evaluating them in the averaged state is not equivalent to averaging their trajectory-resolved values.  The measurement record is consequently part of the physical specification of the prepared ensemble, not bookkeeping that can be discarded without changing the problem.

In this Letter, we show that the class-A adaptive dynamics prepares a wall-localized critical ensemble from a broad class of initial states, including the maximally mixed state.  The interface contribution to the entanglement follows a log-chord form with a coefficient consistent with a $c=1$ chiral fermion, while correlations fitted on controlled intermediate windows are consistent with the corresponding value $\Delta=2$.  Modular-charge packets launched on opposite walls acquire oppositely signed drifts, providing evidence for chirality without assuming a quantized transport coefficient.  Independently, the finite-time tangent spectrum contains a near-marginal mode localized at the wall; we use this as a dynamical diagnostic and do not infer its thermodynamic scaling from the available sizes.  These results connect the boundary of the monitored dynamics to a universal ensemble of states that it prepares and continuously stabilizes, without postselection.

\begin{figure*}[t]
  \figplaceholder{2.1in}{%
  \textbf{FIG.~1 PLACEHOLDER ($2\times2$ panels).}\\[2pt]
  (a) Bilayer adaptive-circuit schematic with programmed mass profile: topological slab $\atopo=1$ for $x\in[6,14]$, trivial exterior $\atriv=30$, walls at $x=6,14$; truncated OW modes ($\nsh$ shells) and on-site fSWAP feedback. [schematic]\\
  (b) Ensemble taxonomy: Born-sampled perfect-correction trajectories $\Gtraj$, postselected branch $\Gps$, mean channel $\Gbar$, Lindblad limit. [schematic]\\
  (c) Exact domain-wall benchmark: strip-resolved entanglement-contour log-chord fits; full-width slope $0.344 \simeq 2\times 0.169$ (per-wall windows). [data: \texttt{exact\_DW\_benchmark\_notes}, \texttt{notebooks/exact\_DW\_GS\_benchmark.ipynb}]\\
  (d) Exact wall correlator: $\overline{|C(y)|^2}$ vs.\ chord distance, slope $-2.02$ on wall columns, steep interior decay. [same source]}
  \caption{Adaptive steering toward a domain-wall Chern insulator.  (a)~Geometry and circuit; (b)~the four dynamical descriptions used in this work; (c,d)~exact ground-state benchmark for the chiral interface criticality: two additive chiral wall contributions to the entanglement contour and a chord-law correlator with exponent~$2$.}
  \label{fig:setup}
\end{figure*}

\emph{Adaptive circuit and domain-wall geometry.---}We consider spinless complex fermions $\hat c_{\bm r,\mu}$ ($\mu=1,2$) on an $N_x\times N_y$ square lattice with periodic boundary conditions, plus an identical ancillary layer.  The target-state structure is generated by the two-band Chern-insulator parent Hamiltonian
\begin{equation}
  h(\bm k)=\sin k_x\,\sigma^x+\sin k_y\,\sigma^y+\big(\alpha-\cos k_x-\cos k_y\big)\sigma^z ,
  \label{eq:parent}
\end{equation}
whose lower band carries Chern number $-1$ for $\alpha\in(0,2)$ and $0$ for $\alpha>2$~\cite{BhuiyanPanJian2026}.  Following Ref.~\cite{BhuiyanPanJian2026}, the circuit measures the number operators $\hat\calN_{\bm r,\nu,\pm}=\hat\chi^\dagger_{\bm r,\nu,\pm}\hat\chi_{\bm r,\nu,\pm}$ of localized overcomplete-Wannier (OW) modes $\hat\chi_{\bm r,\nu,\pm}$ ($\nu=A,B$) built from the band projectors of Eq.~\eqref{eq:parent}; the target state is the unique common eigenstate with all lower-band modes filled and all upper-band modes empty.  After each Born-sampled measurement, an undesired outcome triggers an on-site fermionic SWAP with the ancillary layer that corrects it with finite probability (``perfect correction'')---the feedback target is deterministic, but the measurement outcome is not, so the dynamics remains a genuinely stochastic Markov process on covariance matrices $\Gtraj=2C_\xi-\mathbf{1}$, $[C_\xi]_{ij}=\Tr(\hat\rho_\xi \hat c^\dagger_i\hat c_j)$, labeled by the record $\xi$.  All operations are local; truncating each OW mode to $\nsh$ shells makes them strictly finite-ranged and the circuit depth per cycle $\mathcal{O}(1)$~\cite{BhuiyanPanJian2026}.

The new ingredient is a \emph{programmed interface}, which can be built adaptively in two equivalent ways.  In the first, the mass parameter in Eq.~\eqref{eq:parent} is position dependent: the circuit measures OW occupancies of a topological target ($\alpha(x)=\atopo=1$) inside a slab bounded by two walls (at $x=6,14$ for $N_x=20$) and of a trivial target ($\alpha(x)=\atriv=30$) in the exterior---two distinct target topologies steered on complementary regions, producing a soft wall.  In the second, the measured OW modes are additionally truncated at the sector boundary containing their center; this \texttt{dw\_truncation} implements the hard-wall, infinite-exterior-mass limit of the OW support---an active boundary condition that exposes the interface of a uniformly topological region, not a numerical convenience.  The main text uses the hard-wall construction; the supplement verifies that all results persist for soft walls at finite $\atriv$~\cite{SM}.  In either construction, the finite support of the measured modes introduces spatial disorder at the level of measurement outcomes, so the steady state is an \emph{ensemble} of pure, weakly disordered Gaussian Chern insulators rather than a single wavefunction.

A nonlinear observable $O$ defines three inequivalent objects,
\begin{equation}
  \avg{O[\Gtraj]}\;\neq\;O[\Gbar]\;\neq\;O[\Gps],
  \label{eq:taxonomy}
\end{equation}
where the overline denotes the Born-weighted trajectory average, $\Gbar=\avg{\Gtraj}$ is the mean covariance evolved by the deterministic averaged channel, and $\Gps$ is the deterministic postselected branch in which every outcome is forced to its target value.  Perfect correction realizes the first object; postselection and the mean channel are controls~\cite{JianShapourian2023,SM}.  All trajectory data below are averages of nonlinear observables over Born-sampled corrected trajectories.

To calibrate the interface criticality we use the exact ground state $G=2P_--\mathbf{1}$ of the domain-wall Hamiltonian [Fig.~\ref{fig:setup}(c,d)]: each wall hosts one chiral complex fermion, so a full-width strip of height $A_y$ obeys the log-chord law
\begin{equation}
  \overline{S}(A_y)=m\,\log\!\Big[\tfrac{N_y}{\pi}\sin\!\big(\tfrac{\pi A_y}{N_y}\big)\Big]+b,
  \qquad m_\ast=\frac{c}{3}=\frac{1}{3},
  \label{eq:logchord}
\end{equation}
with the two walls contributing additively ($0.344\simeq2\times0.169$ in the exact benchmark), and the squared correlator along a wall decays with chord exponent $\Delta=2$ (fitted slope $-2.02$)~\cite{SM}.

\begin{figure}[t]
  \figplaceholder{2.9in}{%
  \textbf{FIG.~2 PLACEHOLDER ($2\times2$ panels).}\\[2pt]
  (a) Total entropy vs.\ cycle from maximally mixed initial states: perfect correction purifies to near zero; postselection saturates at $2\log 2$ (dotted line). [data: \texttt{colab\_small\_system\_testing} purification summaries]\\
  (b) Intrinsic charge variance $\Tr C_\xi(1-C_\xi)$ vs.\ cycle, $N_y=30,40,50$: collapse from $\mathcal{O}(N_xN_y)$ to $0.067$--$0.097$ by $t=2N_y$. [data: \path{purification_total_charge_variance_sample_stats_vs_cycle.csv}]\\
  (c) Charge-sharpening crossover vs.\ $\atopo$: threshold $\mathrm{Var}\,Q<10^{-2}$ first crossed at $\atopo\simeq1.9$--$2.0$ (PC) and $1.7$--$1.8$ (PS). [data: \texttt{threshold\_crossing\_summary.csv}]\\
  (d) Centered total charge statistics: percent spread $0.534\%\to0.449\%\to0.413\%$ for $N_y=30,40,50$ ($\propto 1/\sqrt{N}$ guide). [data: \path{purification_centered_total_charge_sample_stats_vs_cycle.csv}]}
  \caption{Purification and charge sharpening of the steered ensemble.  (a)~Protocol contrast from maximally mixed initial states; (b)~intrinsic charge-variance collapse; (c)~sharpening crossover tracking the topological transition of the target state at $\alpha=2$; (d)~sub-extensive filling fluctuations: the filling fraction concentrates toward $1/2$ as $1/\sqrt{N}$.}
  \label{fig:purification}
\end{figure}

\emph{Purification and charge sharpening.---}We first establish that the steering succeeds from the worst-case initial condition.  Starting from the maximally mixed state at $N_x=20$, $N_y=30,40,50$ ($100$ trajectories, $2N_y$ cycles), the intrinsic charge variance $\Tr C_\xi(1-C_\xi)$---which measures the total occupation uncertainty of the Gaussian state---collapses from its maximally mixed value $N_xN_y/2$ to $0.067$, $0.085$, $0.097$ respectively, a suppression of more than four orders of magnitude, while the total entropy purifies in parallel [Fig.~\ref{fig:purification}(a,b)].  Sweeping the slab mass, the sharpening threshold is crossed at $\atopo\simeq1.9$--$2.0$, tracking the topological transition of the target state at $\alpha=2$ [Fig.~\ref{fig:purification}(c)]; this is a feedback-induced sharpening crossover tied to the target-state topology, distinct from the measurement-rate-tuned charge-sharpening transitions of monitored U(1) circuits~\cite{Agrawal2022,Barratt2022}.  Postselection, by contrast, leaves a robust residual entropy $2\log2$ in the same geometries [Fig.~\ref{fig:purification}(a)]: the forced branch is a useful control but not the physical ensemble, and we treat it as such throughout~\cite{SM}.

Crucially, the steered ensemble remains \emph{charge fluctuating}: the sample-to-sample total charge wanders by $6$--$8$ particles (standard deviation) at these sizes [Fig.~\ref{fig:purification}(d)].  These fluctuations do not threaten the chiral universality, for a simple reason: each trajectory is a filled lower band up to occupations of in-gap interface modes, and the effective chiral-edge description depends only on the Fermi point lying on the chiral branch \emph{inside the bulk gap}.  Filling wander merely translates the interface Fermi momentum along the chiral branch, leaving $c=1$ and $\Delta=2$ untouched.  Quantitatively, the wander is small compared with the in-gap budget of ${\sim}2N_y$ wall states, and the filling \emph{fraction} concentrates toward one half with system size: the percent spread falls monotonically as $0.534\%\to0.449\%\to0.413\%$ for $N_y=30,40,50$, consistent with sub-extensive $\sqrt{N}$ charge fluctuations, while the mean offset ($\lesssim0.32\%$) is smaller than the spread and hence statistically consistent with zero~\cite{SM}.  The universality is gap-protected; the filling is not pinned, and need not be.

\begin{figure}[t]
  \figplaceholder{2.9in}{%
  \textbf{FIG.~3 PLACEHOLDER ($2\times2$ panels).}\\[2pt]
  (a) $\overline{S}(A_y)$ vs.\ log-chord coordinate for $N_y=40$--$120$ with $c/3$ guide; $R^2>0.9999$. [data: \texttt{colab\_large\_entanglement\_scaling\_N20} run summaries]\\
  (b) Fitted slope $m$ vs.\ $N_y$: $0.3615, 0.3555, 0.3532, 0.3470$ for $N_y=40$--$100$, flowing toward $1/3$ (dashed). [same source, \texttt{fit\_rows.csv}]\\
  (c) Window-controlled squared-correlator exponent vs.\ $N_y$ with $\Delta=2$ guide: $2.13, 2.12, 2.12$ on $r\in[2,N_y/4]$; greyed default-window values (crossover-dominated). [data: \path{equilibrium_vs_dynamics_selected_window_exponents.csv}]\\
  (d) Modular-charge center-of-mass drift $\Delta y(t)$ for packets launched on opposite walls: opposite signs. [data: \texttt{modular\_charge\_spreading} summaries]}
  \caption{Chiral criticality of the steered interface ensemble.  (a,b)~Log-chord entanglement scaling flowing toward the chiral value $c/3=1/3$; (c)~correlator exponents consistent with $\Delta=2$ on controlled windows; (d)~oppositely signed modular-charge drift on opposite walls---signed-drift evidence of chirality.}
  \label{fig:critical}
\end{figure}

\emph{Chiral criticality of the interface ensemble.---}Figure~\ref{fig:critical} presents the central result.  The trajectory-averaged full-width strip entropy follows the log-chord form of Eq.~\eqref{eq:logchord} with extraordinary fit quality ($R^2>0.9999$) across all sizes, and the fitted slope decreases monotonically toward the chiral target, $m=0.3615,\,0.3555,\,0.3532,\,0.3470$ for $N_y=40,60,80,100$ at depth $t=N_y/2$ [Fig.~\ref{fig:critical}(a,b)].  We quote this as ``near $1/3$ with finite-size and finite-depth corrections'': the residual slope excess $\delta m\simeq0.015$--$0.026$ decreases with both $N_y$ and circuit depth, and dedicated diagnostics attribute it to local wall defects and incomplete late-time convergence rather than to global charge imbalance or averaging-order artifacts~\cite{SM}.  The squared correlator along the wall is a sharper probe but a more delicate fit: on controlled intermediate windows $r\in[2,N_y/4]$ the exponent is $2.13$, $2.12$, $2.12$ for $N_y=30,40,50$, consistent with the chiral value $\Delta=2$, whereas naive long-distance windows are contaminated by finite-size curvature and should not be read as universality data~\cite{SM}.

Logarithmic entanglement alone fixes a critical sector, not its orientation.  To probe chirality directly we evolve localized charge packets under the modular Hamiltonian $h_{\rm mod}=-2\,\mathrm{arctanh}\,G_A$ of the half-cylinder reduction of the prepared covariance.  Packets launched on opposite walls acquire center-of-mass drifts of opposite sign [Fig.~\ref{fig:critical}(d)], as required for counter-propagating chiral interfaces on a periodic geometry; charge conservation in these diagnostics holds to numerical precision.  We present this as signed-drift evidence---velocity magnitudes carry finite-size scatter---and note that it is the dynamical-transport counterpart of static modular-data probes of the chiral central charge~\cite{KimShiKatoAlbert2022,KimShiKatoAlbert2022b,LiHaldane2008,ChenVidal2014}.  Together, Figs.~\ref{fig:critical}(a)--(d) identify the steady-state interface sector as a $c=1$ chiral critical liquid, prepared and persistently re-stabilized by strictly local measurement and feedback.

\emph{The trajectory-averaged skeleton.---}What part of this physics survives if one does not resolve trajectories?  Every trajectory-averaged adaptive circuit maps onto an associated Lindbladian~\cite{BhuiyanPanJian2026,Guo2025}: the Born-averaged dynamics defines a deterministic channel $\Gbar_{n+1}=\Phi(\Gbar_n)$ whose continuous-time embedding carries gain, loss, and dephasing parts~\cite{BhuiyanPanJian2026}.  At this averaged level, Goldstein's no-go theorem applies with full force---a finite-range Lindbladian cannot relax at a finite rate to a unique pure topological state in $d\geq2$~\cite{Goldstein2019}---and indeed the averaged steady state in the domain-wall geometry is mixed.  It is the trajectory-resolved protocol, with finite-support measurements and local feedback, that circumvents the obstruction: purity and chiral criticality live in the individual Gaussian trajectories, not in their average.  The averaged dynamics is moreover \emph{interacting}: the dephasing dissipators carry quartic jump operators $\hat\calN_{\bm r,\nu,\pm}$, so although every individual trajectory is Gaussian, the trajectory-averaged state is generically non-Gaussian~\cite{BhuiyanPanJian2026,Dolgirev2020,BarthelZhang2022}---placing it outside the quadratic-Lindbladian classification~\cite{LieuMcGinleyCooper2020,AltlandFleischhauerDiehl2021}.  Two-point functions nevertheless obey a closed (nonlinear) equation~\cite{BarthelZhang2022,BhuiyanPanJian2026}, and its steady state in the domain-wall geometry retains a striking feature: the $k_y$-resolved occupation spectrum of $\Gbar$ exhibits \emph{edge bands}---branches interpolating between the near-filled and near-empty bulk sectors, localized on the walls [Fig.~\ref{fig:stability}(e)].  This is the interacting, internal-interface counterpart of edge structure known in engineered quadratic dissipators~\cite{Diehl2011,Bardyn2013,BudichZollerDiehl2015,YangMoligniniBergholtz2023,HegdeEhmckeMeng2023}.

What averaging removes is the criticality itself.  For any nonlinear observable the ordering matters [Eq.~\eqref{eq:taxonomy}]; for squared correlators the statement is exact,
\begin{equation}
  \avg{|C_\xi(r,r')|^2}=|\bar C(r,r')|^2+\avg{|C_\xi(r,r')-\bar C(r,r')|^2},
  \label{eq:decomp}
\end{equation}
and the trajectory-fluctuation term---deleted by construction in any mean-channel or Lindblad treatment---carries the power law [Fig.~\ref{fig:stability}(f)].  The same holds for the logarithmic entanglement and the charge-sector statistics.  The dichotomy is sharp: \emph{the wall's spectral skeleton survives averaging; its criticality does not.}  Trajectory resolution---implemented here without postselection, since perfect correction is the physical ensemble---is precisely the resource that converts mixed-state interface structure into pure-state chiral critical matter.

\begin{figure*}[t]
  \figplaceholder{3.0in}{%
  \textbf{FIG.~4 PLACEHOLDER ($3\times2$ panels).}\\[2pt]
  (a) Finite-time Lyapunov spectra of the single-particle transfer frame: uniform vs.\ domain-wall geometry. [data: \texttt{colab\_lyapunov}]\\
  (b) Lyapunov gap $\Dlam$ vs.\ $N_y$ (log scale): uniform $0.72$--$0.77$ (PC), $\sim0.95$ (PS); domain wall $0.027,0.024,0.020$ (PC), decreasing with $N_y$. [data: \texttt{colab\_lyapunov} run index]\\
  (c) Choi particle-transfer gap $\Dp$ vs.\ cycle: $\mathcal{O}(10^{-2})$ in the DW slab vs.\ large transient gaps in the uniform control; censoring window marked. [data: \texttt{colab\_regularized\_choi\_transfer\_matrix}, \texttt{choi\_covariance\_cpu}]\\
  (d) Minimum-gap eigenvector density $|\psi|^2(x,y)$: localized on wall columns; inset: per-sample Chern marker vs.\ $\Dlam$. [data: \texttt{colab\_lyapunov}]\\
  (e) $k_y$-resolved occupation spectrum of the steady-state mean covariance $\Gbar$: wall-localized edge branches interpolating between filled and empty sectors. [to generate: \texttt{classA\_U1FGTN.run\_markov\_channel} / \texttt{CI\_Lindblad\_DW.py}]\\
  (f) $|\bar C(y)|^2$ vs.\ $\avg{|C_\xi(y)|^2}$ along the wall: the trajectory-fluctuation term of Eq.~(\ref{eq:decomp}) carries the power law. [to generate: same entry points + trajectory data]}
  \caption{Stability and the trajectory-averaged skeleton of the interface sector.  (a--d)~Transfer diagnostics: the uniform topological dynamics is strongly contracting while the domain-wall geometry hosts a near-marginal, wall-localized channel; two independent diagnostics (Lyapunov transfer frame and Choi particle transfer) agree qualitatively.  (e,f)~The interacting trajectory-averaged dynamics retains wall-localized occupation-spectrum edge bands but loses the critical correlations, which reside entirely in the trajectory-fluctuation term.}
  \label{fig:stability}
\end{figure*}

\emph{A near-marginal interface channel.---}The criticality of the steered ensemble has a counterpart in the stability spectrum of the dynamics itself.  We linearize the stochastic covariance update along realized trajectories and accumulate finite-time Lyapunov exponents of the single-particle transfer frame (QR-reorthogonalized once per cycle, normalized per cycle), with gap proxy $\Dlam=\min_i|\lambda_i|$~\cite{Ludwig2013,SM}.  In the uniform topological geometry the dynamics is strongly contracting: all exponents are strictly negative with $\Dlam\simeq0.72$--$0.77$ under perfect correction (and $\simeq0.95$ under postselection), plateauing within ${\sim}10$ cycles and showing no size dependence.  In the domain-wall geometry the gap is suppressed by one to two orders of magnitude, $\Dlam\simeq0.027,\,0.024,\,0.020$ for $N_y=30,40,50$, decreases with circumference, and does not plateau within $t=N_y$ [Fig.~\ref{fig:stability}(a,b)].  The minimum-gap tangent vector is concentrated on the wall columns and extended along the wall, and trajectories with smaller gaps carry degraded Chern markers [Fig.~\ref{fig:stability}(d)]---a slow, neutral deformation of the chiral interface rather than a bulk instability.  An independent diagnostic, the finite-sector particle-transfer gap $\Dp$ extracted from the regularized cumulative Choi covariance via $T_{\rm p}=-(\mathbf{1}+\Sigma_{LL})^{-1}\Sigma_{LR}$, shows the same contrast: $\Dp\sim10^{-2}$ in the wall sector against large transient gaps in the uniform control [Fig.~\ref{fig:stability}(c)]~\cite{SM}.  We emphasize that these are distinct linearizations whose qualitative agreement---not operator identity---supports the picture; with three circumferences and depths $t=N_y$, the data establish a near-marginal interface channel but not its thermodynamic scaling form, which we leave open.  Physically, the near-zero channel is the dynamical fingerprint of the gapless interface: the steering contracts onto the critical ensemble quickly in the bulk but only marginally along the wall, which is also visible in the algebraically slow purification of the wall region~\cite{BhuiyanPanJian2026}.  This realizes, within a postselection-free preparation and at internal programmable interfaces, the gapless boundary Lyapunov modes anticipated by transfer-matrix classifications of monitored dynamics~\cite{XiaoKawabata2024,Oshima2025}.

\emph{Discussion.---}These results establish adaptive measurement-feedback circuits as a platform for \emph{measurement-engineered chiral interfaces}: gapless chiral matter created where it is wanted, with the location, number, and confinement of chiral channels programmed through the measurement schedule.  Sweeping the exterior mass confirms the dial: with a topological interior, the extracted interface central charge counts the exposed chiral channels, dropping from $c_{\rm est}\approx2$ to $c_{\rm est}\approx1$ as the exterior crosses its topological transition~\cite{SM}.  Unlike chiral interface transport in materials~\cite{RosenEtAl2017,UpadhyayaTserkovnyak2016} or synthetic-dimension edge states in analog simulators~\cite{Mancini2015,Stuhl2015}, the present route is digital, postselection-free, and self-correcting: the same measurements that prepare the critical ensemble continuously stabilize it, with a finite threshold against coherent noise inherited from the bulk protocol~\cite{BhuiyanPanJian2026}.  Because the dynamics is Gaussian, classical benchmarking is efficient at scale; because feedback is on-site, the protocol maps directly onto midcircuit-measurement hardware via fermion-to-qubit encodings~\cite{BravyiKitaev2002,Altman2021,Bluvstein2024} and is native to emerging fermionic processors~\cite{GonzalezCuadra2023,Ott2025,Schuckert2024}.

From the spacetime perspective of the dynamical bulk--boundary correspondence~\cite{BhuiyanPanJian2026}, the $1{+}1$d wall in $2{+}1$d spacetime carries a chiral critical ensemble---the higher-dimensional generalization of the $0{+}1$d dynamical domain-wall modes of monitored topological chains~\cite{PanShapourianJian2025}, and an answer to the question raised there: the modes are chiral, critical, and near-marginal.  Our criticality is engineered at an interface of gapped area-law dynamics, by design rather than by tuning to a bulk measurement-induced transition~\cite{AlbertonBuchholdDiehl2021,Fava2023,Poboiko2023,ChahineBuchhold2024}.  Several questions remain open: the thermodynamic scaling of the near-marginal channel; an analytic theory of the steered edge, akin to low-energy edge theories in equilibrium; the fate of the dichotomy between the averaged skeleton and trajectory-resolved criticality under interactions; and whether interacting measurements can stabilize chiral critical matter in symmetry classes where Gaussian measurements cannot~\cite{BhuiyanPanJian2026}.

\begin{acknowledgments}
We thank [\,\dots\,] for helpful discussions.  A.B.\ is supported by the Natural Sciences and Engineering Research Council of Canada through the Postgraduate Scholarships---Doctoral program.  H.P.\ is supported by US-ONR Grant No.\ N00014-23-1-2357.  C.-M.J.\ is supported by the Alfred P. Sloan Foundation through a Sloan Research Fellowship.
\end{acknowledgments}

\begin{thebibliography}{99}

\bibitem{CallanHarvey1985}
C.~G. Callan and J.~A. Harvey,
``Anomalies and fermion zero modes on strings and domain walls,''
Nucl. Phys. B \textbf{250}, 427 (1985).

\bibitem{Hatsugai1993}
Y.~Hatsugai,
``Chern number and edge states in the integer quantum Hall effect,''
Phys. Rev. Lett. \textbf{71}, 3697 (1993).

\bibitem{QiZhangRMP}
X.-L. Qi and S.-C. Zhang,
``Topological insulators and superconductors,''
Rev. Mod. Phys. \textbf{83}, 1057 (2011).

\bibitem{HasanKaneRMP}
M.~Z. Hasan and C.~L. Kane,
``Colloquium: Topological insulators,''
Rev. Mod. Phys. \textbf{82}, 3045 (2010).

\bibitem{ChiuTeoSchnyderRyu2016}
C.-K. Chiu, J.~C.~Y. Teo, A.~P. Schnyder, and S.~Ryu,
``Classification of topological quantum matter with symmetries,''
Rev. Mod. Phys. \textbf{88}, 035005 (2016).

\bibitem{Kitagawa2010}
T.~Kitagawa, E.~Berg, M.~Rudner, and E.~Demler,
``Topological characterization of periodically driven quantum systems,''
Phys. Rev. B \textbf{82}, 235114 (2010).

\bibitem{Lindner2011}
N.~H. Lindner, G.~Refael, and V.~Galitski,
``Floquet topological insulator in semiconductor quantum wells,''
Nat. Phys. \textbf{7}, 490 (2011).

\bibitem{Rudner2013}
M.~S. Rudner, N.~H. Lindner, E.~Berg, and M.~Levin,
``Anomalous edge states and the bulk-edge correspondence for periodically driven two-dimensional systems,''
Phys. Rev. X \textbf{3}, 031005 (2013).

\bibitem{Harper2020}
F.~Harper, R.~Roy, M.~S. Rudner, and S.~L. Sondhi,
``Topology and broken symmetry in Floquet systems,''
Annu. Rev. Condens. Matter Phys. \textbf{11}, 345 (2020).

\bibitem{RosenEtAl2017}
I.~T. Rosen \emph{et al.},
``Chiral transport along magnetic domain walls in the quantum anomalous Hall effect,''
npj Quantum Mater. \textbf{2}, 69 (2017).

\bibitem{UpadhyayaTserkovnyak2016}
P.~Upadhyaya and Y.~Tserkovnyak,
``Domain wall in a quantum anomalous Hall insulator as a magnetoelectric piston,''
Phys. Rev. B \textbf{94}, 020411(R) (2016).

\bibitem{Mancini2015}
M.~Mancini \emph{et al.},
``Observation of chiral edge states with neutral fermions in synthetic Hall ribbons,''
Science \textbf{349}, 1510 (2015).

\bibitem{Stuhl2015}
B.~K. Stuhl, H.-I. Lu, L.~M. Aycock, D.~Genkina, and I.~B. Spielman,
``Visualizing edge states with an atomic Bose gas in the quantum Hall regime,''
Science \textbf{349}, 1514 (2015).

\bibitem{Skinner2019}
B.~Skinner, J.~Ruhman, and A.~Nahum,
``Measurement-induced phase transitions in the dynamics of entanglement,''
Phys. Rev. X \textbf{9}, 031009 (2019).

\bibitem{LiChenFisher2018}
Y.~Li, X.~Chen, and M.~P.~A. Fisher,
``Quantum Zeno effect and the many-body entanglement transition,''
Phys. Rev. B \textbf{98}, 205136 (2018).

\bibitem{Fisher2023}
M.~P.~A. Fisher, V.~Khemani, A.~Nahum, and S.~Vijay,
``Random quantum circuits,''
Annu. Rev. Condens. Matter Phys. \textbf{14}, 335 (2023).

\bibitem{GullansHuse2020}
M.~J. Gullans and D.~A. Huse,
``Dynamical purification phase transition induced by quantum measurements,''
Phys. Rev. X \textbf{10}, 041020 (2020).

\bibitem{Roy2020}
S.~Roy, J.~T. Chalker, I.~V. Gornyi, and Y.~Gefen,
``Measurement-induced steering of quantum systems,''
Phys. Rev. Res. \textbf{2}, 033347 (2020).

\bibitem{Puente2024}
D.~A. Puente, F.~Motzoi, T.~Calarco, G.~Morigi, and M.~Rizzi,
``Quantum state preparation via engineered ancilla resetting,''
Quantum \textbf{8}, 1299 (2024).

\bibitem{Smith2024}
K.~C. Smith, A.~Khan, B.~K. Clark, S.~M. Girvin, and T.-C. Wei,
``Constant-depth preparation of matrix product states with adaptive quantum circuits,''
PRX Quantum \textbf{5}, 030344 (2024).

\bibitem{LuLessaKimHsieh2022}
T.-C. Lu, L.~A. Lessa, I.~H. Kim, and T.~H. Hsieh,
``Measurement as a shortcut to long-range entangled quantum matter,''
PRX Quantum \textbf{3}, 040337 (2022).

\bibitem{TVV2023}
N.~Tantivasadakarn, R.~Verresen, and A.~Vishwanath,
``Hierarchy of topological order from finite-depth unitaries, measurement, and feedforward,''
PRX Quantum \textbf{4}, 020339 (2023).

\bibitem{BhuiyanPanJian2026}
A.~Bhuiyan, H.~Pan, and C.-M. Jian,
``Free-fermion dynamics with measurements: Topological classification and adaptive preparation of topological states,''
Phys. Rev. Res. \textbf{8}, 023147 (2026), arXiv:2507.13437.

\bibitem{JianBauerKeselmanLudwig2022}
C.-M. Jian, B.~Bauer, A.~Keselman, and A.~W.~W. Ludwig,
``Criticality and entanglement in nonunitary quantum circuits and tensor networks of noninteracting fermions,''
Phys. Rev. B \textbf{106}, 134206 (2022).

\bibitem{Ludwig2013}
A.~W.~W. Ludwig, H.~Schulz-Baldes, and M.~Stolz,
``Lyapunov spectra for all ten symmetry classes of quasi-one-dimensional disordered systems of non-interacting fermions,''
J. Stat. Phys. \textbf{152}, 275 (2013).

\bibitem{PanShapourianJian2025}
H.~Pan, H.~Shapourian, and C.-M. Jian,
``Topological modes in monitored quantum dynamics,''
Phys. Rev. B \textbf{112}, 144301 (2025).

\bibitem{NahumSkinner2020}
A.~Nahum and B.~Skinner,
``Entanglement and dynamics of diffusion-annihilation processes with Majorana defects,''
Phys. Rev. Res. \textbf{2}, 023288 (2020).

\bibitem{KellsMeidanRomito2023}
G.~Kells, D.~Meidan, and A.~Romito,
``Topological transitions in weakly monitored free fermions,''
SciPost Phys. \textbf{14}, 031 (2023).

\bibitem{XiaoKawabata2024}
Z.~Xiao and K.~Kawabata,
``Topology of monitored quantum dynamics,''
arXiv:2412.06133.

\bibitem{Oshima2025}
H.~Oshima, K.~Mochizuki, R.~Hamazaki, and Y.~Fuji,
``Topology and spectrum in measurement-induced phase transitions,''
Phys. Rev. Lett. \textbf{134}, 240401 (2025).

\bibitem{LiZou2023}
Y.~Li, Y.~Zou, P.~Glorioso, E.~Altman, and M.~P.~A. Fisher,
``Cross entropy benchmark for measurement-induced phase transitions,''
Phys. Rev. Lett. \textbf{130}, 220404 (2023).

\bibitem{GarrattAltman2024}
S.~J. Garratt and E.~Altman,
``Probing postmeasurement entanglement without postselection,''
PRX Quantum \textbf{5}, 030311 (2024).

\bibitem{Guo2025}
J.~Guo, O.~Hart, C.-F. Chen, A.~J. Friedman, and A.~Lucas,
``Designing open quantum systems with known steady states: Davies generators and beyond,''
Quantum \textbf{9}, 1612 (2025).

\bibitem{Diehl2011}
S.~Diehl, E.~Rico, M.~A. Baranov, and P.~Zoller,
``Topology by dissipation in atomic quantum wires,''
Nat. Phys. \textbf{7}, 971 (2011).

\bibitem{BudichZollerDiehl2015}
J.~C. Budich, P.~Zoller, and S.~Diehl,
``Dissipative preparation of Chern insulators,''
Phys. Rev. A \textbf{91}, 042117 (2015).

\bibitem{TTVV2024}
N.~Tantivasadakarn, R.~Thorngren, A.~Vishwanath, and R.~Verresen,
``Long-range entanglement from measuring symmetry-protected topological phases,''
Phys. Rev. X \textbf{14}, 021040 (2024).

\bibitem{Zhu2023}
G.-Y. Zhu, N.~Tantivasadakarn, A.~Vishwanath, S.~Trebst, and R.~Verresen,
``Nishimori's cat: Stable long-range entanglement from finite-depth unitaries and weak measurements,''
Phys. Rev. Lett. \textbf{131}, 200201 (2023).

\bibitem{LeeJiBiFisher2022}
J.~Y. Lee, W.~Ji, Z.~Bi, and M.~P.~A. Fisher,
``Decoding measurement-prepared quantum phases and transitions: from Ising model to gauge theory, and beyond,''
arXiv:2208.11699.

\bibitem{LuZhangVijayHsieh2023}
T.-C. Lu, Z.~Zhang, S.~Vijay, and T.~H. Hsieh,
``Mixed-state long-range order and criticality from measurement and feedback,''
PRX Quantum \textbf{4}, 030318 (2023).

\bibitem{Malz2024}
D.~Malz, G.~Styliaris, Z.-Y. Wei, and J.~I. Cirac,
``Preparation of matrix product states with log-depth quantum circuits,''
Phys. Rev. Lett. \textbf{132}, 040404 (2024).

\bibitem{SahayVerresen2024}
R.~Sahay and R.~Verresen,
``Classifying matrix-product states and projected entangled-pair states preparable via measurement and feedback,''
PRX Quantum \textbf{5}, 040304 (2024).

\bibitem{Friedman2022}
A.~J. Friedman, C.~Yin, Y.~Hong, and A.~Lucas,
``Locality and error correction in quantum dynamics with measurement,''
arXiv:2206.09929.

\bibitem{LuMixedFollowup2025}
Z.~Zhang, T.-C. Lu, S.~Vijay, and T.~H. Hsieh,
``Universal properties of critical mixed states from measurement and feedback,''
arXiv:2503.09597.

\bibitem{ZhuBoundary2023}
Y.~Wu, G.-Y. Zhu, \emph{et al.},
``Triggering boundary phase transitions through bulk measurements in 2D cluster states,''
arXiv:2305.14231.

\bibitem{Dolgirev2020}
P.~E. Dolgirev, J.~Marino, D.~Sels, and E.~Demler,
``Non-Gaussian correlations imprinted by local dephasing in fermionic wires,''
Phys. Rev. B \textbf{102}, 100301(R) (2020).

\bibitem{BarthelZhang2022}
T.~Barthel and Y.~Zhang,
``Solving quasi-free and quadratic Lindblad master equations for open fermionic and bosonic systems,''
J. Stat. Mech. (2022) 113101.

\bibitem{ODea2024}
N.~O'Dea, A.~Morningstar, S.~Gopalakrishnan, and V.~Khemani,
``Entanglement and absorbing-state transitions in interactive quantum dynamics,''
Phys. Rev. B \textbf{109}, L020304 (2024).

\bibitem{Iadecola2023}
T.~Iadecola, S.~Ganeshan, J.~H. Pixley, and J.~H. Wilson,
``Measurement and feedback driven entanglement transition in the probabilistic control of chaos,''
Phys. Rev. Lett. \textbf{131}, 060403 (2023).

\bibitem{Ravindranath2023}
V.~Ravindranath, Y.~Han, Z.-C. Yang, and X.~Chen,
``Entanglement steering in adaptive circuits with feedback,''
Phys. Rev. B \textbf{108}, L041103 (2023).

\bibitem{Piroli2023}
L.~Piroli, Y.~Li, R.~Vasseur, and A.~Nahum,
``Triviality of quantum trajectories close to a directed percolation transition,''
Phys. Rev. B \textbf{107}, 224303 (2023).

\bibitem{SierantTurkeshi2023}
P.~Sierant and X.~Turkeshi,
``Controlling entanglement at absorbing state phase transitions in random circuits,''
Phys. Rev. Lett. \textbf{130}, 120402 (2023).

\bibitem{Bardyn2013}
C.-E. Bardyn, M.~A. Baranov, C.~V. Kraus, E.~Rico, A.~\.{I}mamo\u{g}lu, P.~Zoller, and S.~Diehl,
``Topology by dissipation,''
New J. Phys. \textbf{15}, 085001 (2013).

\bibitem{Goldstein2019}
M.~Goldstein,
``Dissipation-induced topological insulators: A no-go theorem and a recipe,''
SciPost Phys. \textbf{7}, 067 (2019).

\bibitem{BudichDiehl2015}
J.~C. Budich and S.~Diehl,
``Topology of density matrices,''
Phys. Rev. B \textbf{91}, 165140 (2015).

\bibitem{Bardyn2018}
C.-E. Bardyn, L.~Wawer, A.~Altland, M.~Fleischhauer, and S.~Diehl,
``Probing the topology of density matrices,''
Phys. Rev. X \textbf{8}, 011035 (2018).

\bibitem{WawerFleischhauer2021}
L.~Wawer and M.~Fleischhauer,
``Chern number and Berry curvature for Gaussian mixed states of fermions,''
Phys. Rev. B \textbf{104}, 094104 (2021).

\bibitem{YangMoligniniBergholtz2023}
F.~Yang, P.~Molignini, and E.~J. Bergholtz,
``Dissipative boundary state preparation,''
Phys. Rev. Res. \textbf{5}, 043229 (2023).

\bibitem{HegdeEhmckeMeng2023}
S.~S. Hegde, T.~Ehmcke, and T.~Meng,
``Edge-selective extremal damping from topological heritage of dissipative Chern insulators,''
Phys. Rev. Lett. \textbf{131}, 256601 (2023).

\bibitem{LieuMcGinleyCooper2020}
S.~Lieu, M.~McGinley, and N.~R. Cooper,
``Tenfold way for quadratic Lindbladians,''
Phys. Rev. Lett. \textbf{124}, 040401 (2020).

\bibitem{AltlandFleischhauerDiehl2021}
A.~Altland, M.~Fleischhauer, and S.~Diehl,
``Symmetry classes of open fermionic quantum matter,''
Phys. Rev. X \textbf{11}, 021037 (2021).

\bibitem{SM}
See Supplemental Material at [URL] for the circuit algorithm and simulation parameters; the protocol taxonomy and ensemble-ordering derivations; the trajectory-averaged channel and Lindbladian analysis; the exact domain-wall benchmark; the overcomplete-Wannier truncation and form-factor analysis, including the soft-wall ($\mathtt{dw\_truncation}$ off) verification; correlator fit-window and slope-excess diagnostics; modular charge-spreading details; transfer and stability diagnostics; and the circuit-depth bound discussion.

\bibitem{JianShapourian2023}
C.-M. Jian, H.~Shapourian, B.~Bauer, and A.~W.~W. Ludwig,
``Measurement-induced entanglement transitions in quantum circuits of non-interacting fermions: Born-rule versus forced measurements,''
arXiv:2302.09094.

\bibitem{Agrawal2022}
U.~Agrawal, A.~Zabalo, K.~Chen, J.~H. Wilson, A.~C. Potter, J.~H. Pixley, S.~Gopalakrishnan, and R.~Vasseur,
``Entanglement and charge-sharpening transitions in U(1) symmetric monitored quantum circuits,''
Phys. Rev. X \textbf{12}, 041002 (2022).

\bibitem{Barratt2022}
F.~Barratt, U.~Agrawal, S.~Gopalakrishnan, D.~A. Huse, R.~Vasseur, and A.~C. Potter,
``Field theory of charge sharpening in symmetric monitored quantum circuits,''
Phys. Rev. Lett. \textbf{129}, 120604 (2022).

\bibitem{KimShiKatoAlbert2022}
I.~H. Kim, B.~Shi, K.~Kato, and V.~V. Albert,
``Chiral central charge from a single bulk wave function,''
Phys. Rev. Lett. \textbf{128}, 176402 (2022).

\bibitem{KimShiKatoAlbert2022b}
I.~H. Kim, B.~Shi, K.~Kato, and V.~V. Albert,
``Modular commutator in gapped quantum many-body systems,''
Phys. Rev. B \textbf{106}, 075147 (2022).

\bibitem{LiHaldane2008}
H.~Li and F.~D.~M. Haldane,
``Entanglement spectrum as a generalization of entanglement entropy,''
Phys. Rev. Lett. \textbf{101}, 010504 (2008).

\bibitem{ChenVidal2014}
Y.~Chen and G.~Vidal,
``Entanglement contour,''
J. Stat. Mech. (2014) P10011.

\bibitem{BravyiKitaev2002}
S.~B. Bravyi and A.~Y. Kitaev,
``Fermionic quantum computation,''
Ann. Phys. \textbf{298}, 210 (2002).

\bibitem{Altman2021}
E.~Altman \emph{et al.},
``Quantum simulators: Architectures and opportunities,''
PRX Quantum \textbf{2}, 017003 (2021).

\bibitem{Bluvstein2024}
D.~Bluvstein \emph{et al.},
``Logical quantum processor based on reconfigurable atom arrays,''
Nature \textbf{626}, 58 (2024).

\bibitem{GonzalezCuadra2023}
D.~Gonz\'alez-Cuadra \emph{et al.},
``Fermionic quantum processing with programmable neutral atom arrays,''
Proc. Natl. Acad. Sci. USA \textbf{120}, e2304294120 (2023).

\bibitem{Ott2025}
R.~Ott, D.~Gonz\'alez-Cuadra, T.~V. Zache, P.~Zoller, A.~M. Kaufman, and H.~Pichler,
``Error-corrected fermionic quantum processors with neutral atoms,''
Phys. Rev. Lett. \textbf{135}, 090601 (2025).

\bibitem{Schuckert2024}
A.~Schuckert, E.~Crane, A.~V. Gorshkov, M.~Hafezi, and M.~J. Gullans,
``Fermion-qubit fault-tolerant quantum computing,''
arXiv:2411.08955.

\bibitem{AlbertonBuchholdDiehl2021}
O.~Alberton, M.~Buchhold, and S.~Diehl,
``Entanglement transition in a monitored free-fermion chain: From extended criticality to area law,''
Phys. Rev. Lett. \textbf{126}, 170602 (2021).

\bibitem{Fava2023}
M.~Fava, L.~Piroli, T.~Swann, D.~Bernard, and A.~Nahum,
``Nonlinear sigma models for monitored dynamics of free fermions,''
Phys. Rev. X \textbf{13}, 041045 (2023).

\bibitem{Poboiko2023}
I.~Poboiko, P.~P\"opperl, I.~V. Gornyi, and A.~D. Mirlin,
``Theory of free fermions under random projective measurements,''
Phys. Rev. X \textbf{13}, 041046 (2023).

\bibitem{ChahineBuchhold2024}
K.~Chahine and M.~Buchhold,
``Entanglement phases, localization, and multifractality of monitored free fermions in two dimensions,''
Phys. Rev. B \textbf{110}, 054313 (2024).

\end{thebibliography}

\end{document}
